% showcase.tex — adapted teaching deck and preview source
% SPDX-License-Identifier: CC-BY-4.0
%
% Based on Mattia Ippoliti's “Politecnico di Torino Presentation”.
% Adapted and packaged for Apograph by Elia Innocenti. See NOTICE.
\documentclass[aspectratio=169]{beamer}

\usepackage{amsfonts,amsmath,oldgerm}
\usetheme{apographpolito}

\newcommand{\testcolor}[1]{\colorbox{#1}{\textcolor{#1}{test}}~\texttt{#1}}
\newcommand{\hrefcol}[2]{\textcolor{cyan}{\href{#1}{#2}}}
\usefonttheme[onlymath]{serif}

\apographTitleBackground*{theme/assets/background}
\title{Politecnico di Torino Beamer Presentation Theme}
\subtitle{Using \LaTeX\ to prepare slides}
\apographCourse{Master's Degree in Computer Science}
\author{Mattia Ippoliti (original teaching deck)\\Elia Innocenti (Apograph adaptation)}
\date{Apograph showcase}
\apographStudentID{}
\apographEmail{}
\apographSetClosingMessage{Thank you for listening!}

\begin{document}
\maketitle

\begin{frame}{About this adaptation}
This teaching deck is adapted from
\hrefcol{https://www.overleaf.com/latex/templates/politecnico-di-torino-presentation/cnypkdbdyqky}{Mattia Ippoliti's PoliTo presentation},
which in turn derives from
\hrefcol{https://www.overleaf.com/latex/templates/sintef-presentation/jhbhdffczpnx}{Federico Zenith's SINTEF presentation theme}
and related Beamer adaptations.

\vspace{\baselineskip}

Elia Innocenti reorganized the project for Apograph, namespaced its public
commands, separated the minimal starter from this showcase, and added compiling
fallbacks for assets that Apograph cannot redistribute. Full provenance,
modification details, and component licenses are in \texttt{NOTICE}.
\end{frame}

\section{Introduction}

\begin{frame}{Beamer for presentation slides}
\begin{itemize}
\item We assume you can use \LaTeX; if you cannot,
\hrefcol{https://en.wikibooks.org/wiki/LaTeX/}{you can learn it here}
\item Beamer is one of the most popular and powerful document classes for
presentations in \LaTeX
\item Beamer also has a detailed
\hrefcol{https://ctan.org/pkg/beamer}{user manual}
\item Here we present only the most basic features needed to get started
\end{itemize}
\end{frame}

\begin{frame}{Beamer vs. PowerPoint}
Compared to PowerPoint, using \LaTeX\ is useful because:
\begin{itemize}
\item It is not What-You-See-Is-What-You-Get, but
What-You-\emph{Mean}-Is-What-You-Get: you write the content and the computer
does the typesetting
\item It produces a \texttt{pdf}: fewer problems with fonts, formulas, and
program versions
\item It is easier to keep styles, fonts, and highlighting consistent
\item Math typesetting in \TeX\ is excellent:
\begin{equation*}
\mathrm{i}\,\hslash\frac{\partial}{\partial t}\Psi(\mathbf{r},t)=
-\frac{\hslash^2}{2\,m}\nabla^2\Psi(\mathbf{r},t)
+V(\mathbf{r})\Psi(\mathbf{r},t)
\end{equation*}
\end{itemize}
\end{frame}

\begin{frame}[fragile]{Getting Started}
\framesubtitle{Selecting the Apograph PoliTo theme}
Start a presentation with this minimal document:
\begin{block}{Minimum Apograph Beamer Document}
\verb|\documentclass[aspectratio=169]{beamer}|\\
\verb|\usetheme{apographpolito}|\\
\verb|\begin{document}|\\
\verb|\begin{frame}{Hello, world!}|\\
\verb|\end{frame}|\\
\verb|\end{document}|\\
\end{block}
The downloadable project already contains the theme and its dependencies.
\end{frame}

\begin{frame}[fragile]{Title page}
Set the ordinary Beamer metadata and the optional course label in the preamble:
\begin{block}{Commands for the Title Page}
\begin{verbatim}
\title{Sample Title}
\subtitle{Sample subtitle}
\author{First Author, Second Author}
\date{\today}
\apographCourse{Course or event}
\end{verbatim}
\end{block}
Write the title page with \verb|\maketitle|. Use
\verb|\apographTitleBackground{...}| for a full background, or its starred
form for a split view on the right. Missing optional images produce a safe
plain or neutral fallback.
\end{frame}

\begin{frame}[fragile]{Writing a Simple Slide}
\framesubtitle{It's really easy!}
\begin{itemize}[<+->]
\item A typical slide has a bulleted list
\item Items can be uncovered in sequence
\end{itemize}
\begin{block}{Code for a Page with an Itemised List}<+->
\begin{verbatim}
\begin{frame}{Writing a Simple Slide}
  \framesubtitle{It's really easy!}
  \begin{itemize}[<+->]
    \item A typical slide has a bulleted list
    \item Items can be uncovered in sequence
  \end{itemize}\end{frame}
\end{verbatim}
\end{block}
\end{frame}

\section{Personalization}

\apographFootlineColor{apographpolito-yellow}
\begin{frame}[fragile]{Changing Slide Style}
\begin{itemize}
\item Select the white or main-color slide style in the preamble with
\verb|\apographThemeColor{white}| (default) or
\verb|\apographThemeColor{main}|
  \begin{itemize}
  \item Avoid changing it inside the document; Beamer color state is global
  \end{itemize}
\item Change the footline color before a frame with
\verb|\apographFootlineColor{<color>}|
  \begin{itemize}
  \item Theme colors include \testcolor{apographpolito-primary},
  \testcolor{apographpolito-yellow}, \testcolor{apographpolito-green}, and
  \testcolor{apographpolito-darkgreen}
  \item Restore the transparent footline with
  \verb|\apographFootlineColor{}|
  \end{itemize}
\end{itemize}
\end{frame}

\begin{frame}[fragile]{Blocks}
\begin{columns}
\begin{column}{0.3\textwidth}
\begin{block}{Standard Blocks}
These coordinate with the selected theme and footline color.
\begin{verbatim}
\begin{block}{title}
content...
\end{block}
\end{verbatim}
\end{block}
\end{column}
\begin{column}{0.7\textwidth}
\begin{apographcolorblock}[black]{apographpolito-lightgreen}{Colour Blocks}
Choose the block color explicitly. Text is white by default, but an optional
first argument changes it.
\small
\begin{verbatim}
\begin{apographcolorblock}[black]
  {apographpolito-lightgreen}{title}
content...
\end{apographcolorblock}
\end{verbatim}
\end{apographcolorblock}
Other palette choices include \testcolor{apographpolito-purple},
\testcolor{apographpolito-primary}, \testcolor{apographpolito-darkgreen}, and
\testcolor{apographpolito-yellow}.
\end{column}
\end{columns}
\end{frame}

\apographFootlineColor{}
\begin{frame}[fragile]{Using Colours}
\begin{itemize}[<alert@2>]
\item Use \verb|\textcolor{<color name>}{text}| for explicit color
\item The palette exposes stable Apograph names:
  \begin{itemize}
  \item Primary: \testcolor{apographpolito-primary} and its light companion
  \testcolor{apographpolito-light}
  \item Greens: \testcolor{apographpolito-lightgreen},
  \testcolor{apographpolito-green}, and
  \testcolor{apographpolito-darkgreen}
  \item Accents: \testcolor{apographpolito-yellow},
  \testcolor{apographpolito-red}, and \testcolor{apographpolito-purple}
  \end{itemize}
\item Do not abuse colors: \verb|\emph{}| is usually enough
\item Use \verb|\alert{}| to bring the \alert<2->{focus} somewhere
\item<2-|alert@2> If you highlight too much, you do not highlight at all!
\end{itemize}
\end{frame}

\begin{frame}[fragile]{Adding images}
\begin{columns}
\begin{column}{0.7\textwidth}
Adding images works like in normal \LaTeX:
\begin{block}{Code for Adding Images}
\begin{verbatim}
\usepackage{graphicx}
% ...
\includegraphics[width=\textwidth]
  {theme/assets/logo_RGB}
\end{verbatim}
\end{block}
\end{column}
\begin{column}{0.3\textwidth}
\IfFileExists{theme/assets/logo_RGB.png}{%
  \includegraphics[width=\textwidth]{theme/assets/logo_RGB}%
}{%
  \centering\fbox{\parbox[c][18mm][c]{0.8\linewidth}{\centering Optional\\user-provided logo}}%
}
\end{column}
\end{columns}
\end{frame}

\begin{frame}[fragile]{Splitting in Columns}
Splitting a page is easy and common; typically one side has a picture and the
other text:
\begin{columns}
\begin{column}{0.6\textwidth}
This is the first column.
\end{column}
\begin{column}{0.3\textwidth}
And this is the second.
\end{column}
\end{columns}
\begin{block}{Column Code}
\begin{verbatim}
\begin{columns}
  \begin{column}{0.6\textwidth}
    This is the first column
  \end{column}
  \begin{column}{0.3\textwidth}
    And this is the second
  \end{column}
\end{columns}
\end{verbatim}
\end{block}
\end{frame}

\begin{apographchapter}[theme/assets/background_negative]{}{Special Slides}
\begin{itemize}
\item Chapter slides
\item Side-picture slides
\end{itemize}
\end{apographchapter}

\apographFootlineColor{apographpolito-red}
\begin{frame}[fragile]{Chapter slides}
\begin{itemize}
\item Chapter slides behave like frames with a richer layout
\item Open one with
\verb|\begin{apographchapter}[<image>]{<color>}{<title>}|
\item The image is optional; color and title are mandatory arguments
\item The palette includes \testcolor{apographpolito-primary},
\testcolor{apographpolito-darkgreen}, \testcolor{apographpolito-green},
\testcolor{apographpolito-lightgreen}, \testcolor{apographpolito-red},
\testcolor{apographpolito-yellow}, and \testcolor{apographpolito-purple}
\item Matching the following footline to the chapter color can provide useful
visual continuity
\end{itemize}
\end{frame}

\begin{apographsidepic}{theme/assets/background_alternative}{Side-Picture Slides}
\begin{itemize}
\item Open one with
\texttt{\textbackslash begin\{apographsidepic\}\{<image>\}\{<title>\}}
\item Content behaves like an ordinary frame in the left-hand area
\item A neutral placeholder appears when the optional image is absent
\end{itemize}
\end{apographsidepic}

\apographFootlineColor{}
\begin{frame}{Fonts}
\begin{itemize}
\item The paramount task of fonts is being readable
\item There are good choices:
  \begin{itemize}
  \item {\textrm{Use serif fonts only with high-definition projectors}}
  \item {\textsf{Use sans-serif fonts otherwise, or if you prefer them}}
  \end{itemize}
\item And choices that should be used sparingly:
  \begin{itemize}
  \item {\texttt{Avoid monospace for normal prose}}
  \item {\frakfamily Gothic, calligraphic, or decorative fonts should remain
  exceptional}
  \end{itemize}
\end{itemize}
\end{frame}

\begin{frame}[fragile]{Finishing the presentation}
\begin{itemize}
\item Insert a final slide with \verb|\apographBackmatter|
  \begin{itemize}
  \item Change its message with \verb|\apographSetClosingMessage{...}|
  \item Remove the title from it with
  \verb|\apographBackmatter[notitle]|
  \item Set footline text explicitly with
  \verb|\apographFootlineText{...}|
  \end{itemize}
\item The starter uses a 16:9 aspect ratio. Converting a finished presentation
to 4:3 can disturb spacing, so choose the target ratio before authoring.
\end{itemize}
\end{frame}

\section{Summary}

\begin{frame}{Good Luck!}
\begin{itemize}
\item This is enough for an introduction to the theme and Beamer workflow
\item Start real work in \texttt{main.tex}, \texttt{config.tex}, and
\texttt{content/slides.tex}; this showcase is only a reference
\item Corrections and reproducible problems are welcome through
\hrefcol{https://github.com/eliainnocenti/apograph/issues}{Apograph Issues}
\end{itemize}
\end{frame}

\apographBackmatter
\end{document}
